% Part: first-order-logic
% Chapter: beyond
% Section: many-sorted-logic

\documentclass[../../../include/open-logic-section]{subfiles}

\begin{document}

\olfileid{fol}{byd}{msl}

\olsection{బహు-రక తర్కం}

మొదటిస్థాయి తర్కంలో చరాలు, పరిమాణకాలు ఒకే
\tetoken{వ్యక్తి క్షేత్రంపై}{!!{domain}} విస్తరిస్తాయి. అయితే అనేక (పరస్పరం వియుక్తమైన)
\tetoken{వ్యక్తి క్షేత్రాలు}{!!{domain}s} ఉండటం తరచుగా ఉపయోగకరం: ఉదాహరణకు, సంఖ్యల కోసం
\tetoken{ఒక వ్యక్తి క్షేత్రం}{!!a{domain}}, జ్యామితీయ వస్తువుల కోసం \tetoken{ఒక వ్యక్తి క్షేత్రం}{!!a{domain}}, సంఖ్యల నుంచి
సంఖ్యలకు ప్రమేయాల కోసం \tetoken{ఒక వ్యక్తి క్షేత్రం}{!!a{domain}}, అబేలియన్ సమూహాల కోసం
\tetoken{ఒక వ్యక్తి క్షేత్రం}{!!a{domain}}, ఇలా మరెన్నో కావచ్చు.

బహు-రక తర్కం ఇటువంటి చట్రాన్ని అందిస్తుంది. మొదట ``రకాల'' జాబితాను
తీసుకుంటాం---ఒక వస్తువు యొక్క ``రకం'', అది ఏ ``\tetoken{వ్యక్తి క్షేత్రం}{!!{domain}}''లో
ఉండాలో సూచిస్తుంది. తరువాత ప్రతి రకానికీ \tetoken{చరాలు}{!!{variable}s},
పరిమాణకాలు, (సాధారణంగా) ఒక \tetoken{తాదాత్మ్య విధేయం}{!!{identity}} ఉంటాయి. ప్రమేయాలు,
సంబంధాలు తీసుకోగల వస్తువుల రకాల ప్రకారం వాటికీ ``టైపులు''
నిర్ణయిస్తాం. మిగతా విషయాల్లో మొదటిస్థాయి తర్కపు సాధారణ నియమాలనే
ఉంచి, ప్రతి రకానికీ పరిమాణక నియమాల రూపాలను పునరావృతం చేస్తాం.

ఉదాహరణకు, అంతర్జాతీయ సంబంధాలను అధ్యయనం చేయడానికి ఫ్రెంచ్ పౌరులు,
జర్మన్ పౌరులు అనే రెండు రకాల వస్తువులున్న భాషను ఎంచుకోవచ్చు. మొదటి
రకపు వస్తువులకు ``వైన్ తాగుతాడు'' అనే ఏకస్థాన సంబంధం; రెండవ రకపు
వస్తువులకు ``వుర్‌స్ట్ తింటాడు'' అనే మరో ఏకస్థాన సంబంధం; మొదటి
ప్రాచలం మొదటి రకానికీ, రెండవ ప్రాచలం రెండవ రకానికీ చెందిన రెండు
ప్రాచలాలను తీసుకునే ``బహుళజాతి వివాహిత జంటగా ఏర్పడతారు'' అనే
ద్విస్థాన సంబంధం ఉండవచ్చు. ఫ్రెంచ్ పౌరులపై విస్తరించడానికి $a$, $b$,
$c$ చరాలను, జర్మన్ పౌరులపై విస్తరించడానికి $x$, $y$, $z$ చరాలను
వాడితే,
\[
\lforall[a][\lforall[x][(\Atom{\Obj{MarriedTo}}{a,x} \lif
(\Atom{\Obj{DrinksWine}}{a} \lor \lnot \Atom{\Obj{EatsWurst}}{x})]]
\]
అంటే, ఏ ఫ్రెంచ్ వ్యక్తి అయినా ఒక జర్మన్ వ్యక్తిని వివాహం చేసుకొని
ఉంటే, ఆ ఫ్రెంచ్ వ్యక్తి వైన్ తాగుతాడు లేదా ఆ జర్మన్ వ్యక్తి వుర్‌స్ట్
తినడు అని ఇది ప్రకటిస్తుంది.
\sourcecorrection{OLTEFOLBYD-001}{ఆంగ్ల మూలంలోని లోపభూయిష్టమైన \texttt{\string\lforall x]}ను, వెలుపలి పరిమాణకపు ఐచ్ఛిక ప్రాచలాన్ని ముందే మూసివేసే ముద్రణగా చదివి, సమాంతర సంకేతనానికి అనుగుణంగా \texttt{\string\lforall[x][\dots]}గా సరిచేశాం.}

బహు-రక \tetoken{వ్యక్తి క్షేత్రాలలోని}{!!{domain}s} వస్తువులన్నింటినీ ఒకే మొదటిస్థాయి
\tetoken{వ్యక్తి క్షేత్రంలో}{!!{domain}} చేర్చి, రకాలను గుర్తించడానికి ఏకస్థాన \tetoken{విధేయ సంకేతాలను}{!!{predicate}s}
వాడి, పరిమాణకాలను సాపేక్షీకరించడం ద్వారా బహు-రక తర్కాన్ని సహజంగా
మొదటిస్థాయి తర్కంలో ఇమిడ్చవచ్చు. ఉదాహరణకు, పైనున్న ఉదాహరణకు
అనుగుణమైన మొదటిస్థాయి భాషలో, వివరించిన ఇతర సంబంధాలతో పాటు రకపు
నిబంధనలను తొలగించి, ``$\Obj{German}$'', ``$\Obj{French}$'' అనే
ఏకస్థాన \tetoken{విధేయ సంకేతాలు}{!!{predicate}s} ఉంటాయి. $x$ జర్మన్ రకానికి చెందిన
\tetoken{ఒక చరం}{!!a{variable}} అయినప్పుడు, రక-పరిమాణకం $\lforall[x][!A]$ కిందిదిగా
అనువదించబడుతుంది:
\[
\lforall[x][(\Atom{\Obj{German}}{x} \lif !A)].
\]
రకాలు వేరుగా ఉన్నాయని నిర్ధారించే స్వీకృతాలను కూడా చేర్చాలి---ఉదాహరణకు,
$\lforall[x][\lnot (\Atom{\Obj{German}}{x} \land
  \Atom{\Obj{French}}{x})]$---అలాగే ``వైన్ తాగుతాడు'' అనేది
$\Atom{\Obj{French}}{x}$ అనే విధేయాన్ని సంతృప్తిపరచే వస్తువులకు
మాత్రమే వర్తిస్తుందని హామీ ఇచ్చే స్వీకృతాలు మొదలైనవి చేర్చాలి. ఈ
సంప్రదాయాలు, స్వీకృతాలతో బహు-రక \tetoken{వాక్యాలు}{!!{sentence}s} మొదటిస్థాయి
\tetoken{వాక్యాలుగా}{!!{sentence}s}, బహు-రక \tetoken{వ్యుత్పత్తులు}{!!{derivation}s} మొదటిస్థాయి
\tetoken{వ్యుత్పత్తులుగా}{!!{derivation}s} అనువదించబడతాయని చూపడం కష్టం కాదు. అలాగే బహు-రక
\tetoken{నిర్మాణాలు}{!!{structure}s} అనుగుణమైన మొదటిస్థాయి \tetoken{నిర్మాణాలుగా}{!!{structure}s}
``అనువదించబడతాయి'', తిరుగు దిశలోనూ అలాగే జరుగుతుంది. కాబట్టి
బహు-రక తర్కానికి కూడా ఒక సంపూర్ణతా సిద్ధాంతం ఉంటుంది.

\end{document}
